\chapter*{Preface}
\pdfbookmark[0]{Preface}{review.preface}
\label{review:preface}
\markboth{Open Source Evaluatology}{Preface}
A contribution to open source rarely ends where it is recorded. A patch may enter a library used by other projects. A careful review may prevent a failure that never appears in an issue tracker. A maintainer may help a newcomer whose most consequential work appears years later, elsewhere. These activities leave unequal traces. Their effects travel through technical dependencies and human relationships, and their value depends on who experiences them.

Evaluation enters this setting whenever someone must decide what deserves recognition, adoption, attention, or support. A contributor wants their work to be understood. A maintainer must distribute limited review time. An organization needs to understand the components on which it depends. A foundation must decide where assistance can sustain a shared capability. Each question calls for evidence, but the same observation does not carry the same meaning in every decision.

\emph{Open Source Evaluatology: Theory, Methods, and Empirical Studies} develops a foundation for investigating these questions. We establish Open Source Evaluatology as a domain of Evaluatology, grounded in open collaboration and connected to its practice. The parent discipline supplies concepts for examining objects, conditions, references, effects, and value. Open source gives these concepts a setting in which production crosses organizational boundaries, artifacts are widely reused, and benefits and maintenance obligations may fall on different parties.

The book grows out of a sustained programme of research. Activity studies made forms of participation visible and examined the distribution of collaborative work. OpenRank developed a relational account of influence and contribution. OpenDigger made observation, classification, and computation reusable. Contribution-recognition deployments brought evaluation into community practice, where its consequences could themselves be studied. Two doctoral dissertations provide much of the extended theoretical, technical, and empirical work behind these contributions. Our position paper connected this programme to Evaluatology; this monograph develops that connection into a systematic account.

The task is to explain how these achievements belong together. A descriptive measure needs an account of what it represents. A relational model needs a reason for its graph construction. An evaluation used in a course, community, or institution needs evidence about the purposes it serves and the consequences it creates. Tool development makes these inquiries practicable, while shared benchmarks and interdisciplinary research help examine the claims that follow. The book brings these responsibilities into one field without treating them as interchangeable.

Theory and practice work in both directions. Evaluatology guides the choice of object, conditions, comparison, and value. Empirical findings can require those choices to be reconsidered. Neglected coordination work changes the account of contribution. A maintained successor to an archived project changes the boundary of continuity. A public ranking can change the participation it represents. An institutional report can reveal a question that the available indicators cannot answer. Such findings contribute to theory as well as to an immediate application.

The eleven chapters develop this argument in four parts. Part I, \emph{Theory}, establishes the inquiry. Chapter 1 defines the domain and its foundations in Evaluatology. Chapter 2 examines collective capability, stakeholder effects, sustainability, and the consequences of evaluation. Chapter 3 turns these concepts into an evaluation-system design and relates it to completed studies.

Part II, \emph{Methods}, develops the means of inquiry. Chapter 4 examines Activity, from collaboration records and definitions to aggregation and empirical patterns. Chapter 5 develops OpenRank's global and community constructions, their validation, and the interpretation of relational scores. Chapter 6 explains OpenDigger and follows scalable computation and contribution recognition into use. Experiments and deployments establish the methods' capabilities and their limits.

Part III, \emph{Empirical Studies}, examines how the question changes with the setting. Chapter 7 studies contribution recognition through educational deployments. Chapter 8 investigates community sustainability through archival, maintenance, exit, transfer, and renewal. Chapter 9 examines ecosystem governance through Redis and Valkey, organizational contributions, and attribution. Together, the cases connect judgments about a person's work, a continuing capability, and an institutional arrangement.

Part IV, \emph{Applications and Disciplinary Development}, connects these findings to institutions and shared research practice. Chapter 10 develops open source think-tank applications through the global development report, observatories, and OpenDipper. Chapter 11 organizes the disciplinary framework, the open tool ecosystem, OpenPerf's benchmark comparisons, and the field's interdisciplinary foundations. It closes with research questions about comparability, visibility, reactivity, and revision, each grounded in the preceding evidence.

Three distinctions remain central throughout: recorded activity, relational influence, and stakeholder value; the evaluation of a measure and the evaluation of an arrangement that uses it; and reproducible computation and a justified interpretation. Each chapter develops its own question and references. Published results retain their populations, periods, units, and values. Current formulations guide the methodological exposition, while historical findings remain attached to the definitions that produced them. Proposed extensions are distinguished from completed studies.

We write for researchers and graduate students, and for maintainers, community leaders, educators, and institutional practitioners who need to examine the basis of an evaluation. Readers entering through a method can return to its assumptions in Part I; those entering through an application can trace its evidence to Part II. Open source software is the empirical centre of the book. Other forms of open collaboration require their own examination of production, rights, and affected parties.

Our aim is to make collaborative work intelligible, connect evidence to consequential choices, and learn from the consequences of those choices. The resulting foundation is intended to be used, examined, and revised.

\par\vspace{8pt}
{\raggedleft\sffamily\fontsize{10}{13}\selectfont\BookAuthors\par\BookDate\par}
\clearpage
